% ==============================================================================
% APÊNDICES DA DISSERTAÇÃO (PROPECAUT / IFES)
% ==============================================================================

\newpage
\section*{APÊNDICES}
\addcontentsline{toc}{section}{APÊNDICES}

% ------------------------------------------------------------------------------
% APÊNDICE A
% ------------------------------------------------------------------------------
\subsection*{APÊNDICE A -- Catálogo Paramétrico Completo dos Modelos Elétricos de Potência}
\addcontentsline{toc}{subsection}{APÊNDICE A -- Catálogo Paramétrico Completo dos Modelos Elétricos de Potência}

Este apêndice reúne o conjunto exaustivo de propriedades eletrofísicas, térmicas e dimensionais empregadas na parametrização dos modelos do barramento HVDC de 400~V e do gerador fotovoltaico multijunção AM0 desenvolvidos no Capítulo 3, discriminando os valores do nó satelital unitário representativo de 50~kW e da infraestrutura agregada da constelação (1,0~MW útil).

\begin{table}[!ht]
\centering
\footnotesize
\renewcommand{\arraystretch}{0.80}
\caption{Parâmetros detalhados de simulação do barramento HVDC de 400~V e componentes associados.}
\label{tab:apendice_parametros_barramento}
\resizebox{0.95\textwidth}{!}{%
\begin{tabular}{lcccp{6.0cm}}
\hline
\textbf{Parâmetro / Propriedade} & \textbf{Símbolo} & \textbf{Valor Adotado} & \textbf{Unidade} & \textbf{Descrição e Referência Técnica} \\ \hline
Tensão Nominal de Barramento & $V_{nom}$ & 400,0 & V & NASA-HDBK-4007A (NASA, 2026); classe HVDC \\
Tensão de Corte por Subtensão & $V_{UV}$ & 324,0 & V & 3,0 V/célula Li-ion NMC (descarga máxima 100\%) \\
Tensão Máxima de Sobretensão & $V_{OV}$ & 442,8 & V & 4,1 V/célula Li-ion NMC (teto de recarga segura) \\
Potência Máxima Útil de IA (Nó) & $P_{IT,max}$ & 40000 & W & Aceleradores de IA em carga plena no nó (2 racks de 20 kW) \\
Potência em Repouso do Nó de IA & $P_{IT,idle}$ & 2000 & W & Consumo de polarização e standby no eclipse \\
Potência Nominal Bruta do Satélite & $P_{no,nom}$ & 50000 & W & 40 kW IA + 10 kW cargas auxiliares e recarga \\
Eficiência Nominal dos Conversores & $\eta_{conv,nom}$ & 97,5 & \% & Topologia GaN com retificação síncrona a 150 kHz \\
Eficiência Conservadora de Projeto & $\eta_{conv,cons}$ & 90,0 & \% & Pior caso para envelope de dissipação térmica \\
Capacidade do Banco BESS por Nó & $E_{bat,no}$ & 9,0 & kWh & Bateria 108S1P Li-ion NMC ($22{,}5\text{ Ah}$, 60 kg) (NASA, 2021) \\
Capacidade Agregada da Constelação & $E_{bat,tot}$ & 180,0 & kWh & 20 nós satelitais operacionais ($20 \times 9{,}0\text{ kWh}$) \\
Resistência Interna BESS BOL ($20^\circ\text{C}$) & $R_{0,nom}$ & 45 & m$\Omega$ & Medição de impedância em acumuladores NMC de 108S \\
Resistência Interna BESS EOL ($5^\circ\text{C}$) & $R_{0,eol}$ & 85 & m$\Omega$ & Degradação eletroquímica após 10 anos de ciclagem \\
Material dos Cabos de Distribuição & -- & Al 6061-T6 & -- & Densidade $2{,}70\text{ g/cm}^3$, resistividade $2{,}82 \times 10^{-8}\ \Omega\cdot\text{m}$ \\
Bitola dos Condutores Principais & AWG & AWG 4 & -- & Área de seção transversal condutora de $21{,}15\text{ mm}^2$ \\
Comprimento do Chicote Interno & $L_{cabo}$ & 4,5 & m & Distância entre bay de baterias e rack computacional \\
Resistência Ôhmica Nominal dos Cabos & $R_{cab}$ & 14,0 & m$\Omega$ & Condução em alumínio a $20^\circ\text{C}$ (ida e volta 9,0 m) \\
Resistência de Contato dos Terminais & $R_{cont}$ & 2,0 & m$\Omega$ & Terminais crimpados banhados a níquel e prata \\
Indutância Parasita de Distribuição & $L_{dist}$ & 1,8 & $\mu$H & Arranjo blindado com retorno adjacente equilibrado \\
Indutância Desfavorável de Malha & $L_{dist,max}$ & 4,0 & $\mu$H & Pior caso de afastamento e laços de chicote no chassi \\
Capacitância de Filtro do Barramento & $C_{bus}$ & 4,7 & mF & Banco de capacitores cerâmicos/tântalo espaciais \\
Capacitância Degradada em EOL & $C_{bus,eol}$ & 1,5 & mF & Limite inferior considerando perda dielétrica e estresse \\
Resistência de Curto-Circuito Franco & $R_{sc}$ & 1,0 & m$\Omega$ & Estresse extremo de arco condutivo no vácuo \\
Tempo de Abertura Nominal do SSPC & $t_{trip,nom}$ & 15 & $\mu$s & Driver de estado sólido com SiC MOSFET integrado \\
Tempo de Abertura Limite WCA & $t_{trip,max}$ & 100 & $\mu$s & Degradação severa da malha de sensoriamento \\
Limite Térmico do SiC MOSFET & $I^2t_{lim}$ & 1500 & A$^2$s & Limite de queima adiabática de junção (ECSS) \\
Capacitância Térmica do Chassi & $C_{th,elec}$ & 18000 & J/K & Massa estrutural de alumínio de $\approx 20\text{ kg}$ \\
Condutância Térmica ao Loop Heat Pipe & $G_{th}$ & 120 & W/K & Interface condutiva acoplada ao evaporador bifásico \\
\hline
\end{tabular}%
}
\end{table}

\begin{table}[!ht]
\centering
\footnotesize
\renewcommand{\arraystretch}{0.80}
\caption{Parâmetros eletrofísicos e coeficientes térmicos da célula fotovoltaica multijunção AM0.}
\label{tab:apendice_parametros_pv}
\resizebox{0.95\textwidth}{!}{%
\begin{tabular}{lcccp{6.0cm}}
\hline
\textbf{Parâmetro / Propriedade} & \textbf{Símbolo} & \textbf{Valor Adotado} & \textbf{Unidade} & \textbf{Descrição e Referência Técnica} \\ \hline
Constante Solar Extraterrestre AM0 & $G_{AM0}$ & 1361,0 & W/m$^2$ & Wertz et al. (2011); padrão extraterrestre LEO \\
Tecnologia Semicondutora & -- & InGaP/InGaAs/Ge & -- & Tripla junção monolítica III-V (Spectrolab XTJ Prime) \\
Área Ativa adotada no modelo & $A_{cel}$ & 30,18 & cm$^2$ & Geometria escolhida para o cálculo [PROVISÓRIA; não universal do catálogo] \\
Eficiência MPP BOL declarada no datasheet (28°C) & $\eta_{mp}$ & 30,7 & \% & Valor de catálogo; discrepante de $V_{mp}J_{mp}/G$ arredondados \\
Tensão de Circuito Aberto por Célula & $V_{oc,ref}$ & 2,720 & V & Valor de catálogo Spectrolab a $28^\circ\text{C}$ \\
Densidade de Curto-Circuito & $J_{sc,ref}$ & 18,00 & mA/cm$^2$ & Valor de catálogo; corrente depende da área adotada \\
Tensão no Ponto de Máxima Potência & $V_{mp,ref}$ & 2,406 & V & Valor de catálogo Spectrolab a $28^\circ\text{C}$ \\
Densidade no Ponto de Máxima Potência & $J_{mp,ref}$ & 17,50 & mA/cm$^2$ & Valor de catálogo; corrente depende da área adotada \\
Corrente de MPP para $A_{cel}=30{,}18$ cm$^2$ & $I_{mp,ref}$ & 0,52815 & A & Calculada de $J_{mp,ref}A_{cel}$; área [PROVISÓRIA] \\
Potência MPP calculada por célula & $P_{mp,cel}$ & 1,27073 & W & $V_{mp,ref}I_{mp,ref}$; antes de perdas \\
Coeficiente Térmico de $V_{oc}$ & $K_V$ & $-5,60$ & mV/K & Valor de catálogo usado no script \\
Coeficiente Térmico de $J_{sc}$ & $dJ_{sc}/dT$ & $+10,0$ & $\mu$A/(cm$^2$ K) & Valor de catálogo usado no script \\
Coeficiente Térmico de $V_{mp}$ & $dV_{mp}/dT$ & $-6,30$ & mV/K & Valor de catálogo usado nas curvas [VERIFICÁVEL] \\
Coeficiente Térmico de $J_{mp}$ & $dJ_{mp}/dT$ & $+5,0$ & $\mu$A/(cm$^2$ K) & Valor de catálogo usado nas curvas [VERIFICÁVEL] \\
Fator de Idealidade Equivalente & $a_{eq}$ & 3,00 & -- & Parâmetro do modelo concentrado [PROVISÓRIO] \\
Resistência Série Concentrada & $R_s$ & 0,080 & $\Omega$ & Parâmetro do modelo concentrado [PROVISÓRIO] \\
Resistência Shunt de Fuga & $R_{sh}$ & 1200,0 & $\Omega$ & Parâmetro do modelo concentrado [PROVISÓRIO] \\
Corrente de Saturação Equivalente (28°C) & $I_{0,ref}$ & $3{,}6309 \times 10^{-16}$ & A & Recalculada no modelo concentrado [PROVISÓRIO] \\
Número de Células em Série por String & $N_s$ & 166 & -- & Arranjo adotado; $V_{mp}=399{,}4$ V no ponto de catálogo \\
Número de Strings em Paralelo (Nó) & $N_{p,no}$ & 244 & -- & Arranjo adotado com margem de potência instalada \\
Número Total de Células por Satélite & $N_{cel,no}$ & 40504 & -- & $166\times244$; potência MPP ideal de 51,47 kW \\
Área Ativa de Células por Satélite & $A_{ativa,no}$ & 122,24 & m$^2$ & Calculada com área por célula [PROVISÓRIA] \\
Potência MPP BOL por Satélite & $P_{BOL,no}$ & 51,47 & kW & Antes de perdas de conversão e ambientais \\
Strings em Paralelo na Constelação & $N_{p,tot}$ & 4880 & -- & 20 nós operacionais ($20\times244$) \\
Células Totais na Constelação & $N_{total}$ & 810080 & -- & 20 satélites operacionais ($20\times40504$) \\
Área Ativa Total na Constelação & $A_{ativa,tot}$ & 2444,82 & m$^2$ & Calculada; área por célula adotada [PROVISÓRIA] \\
Potência MPP BOL agregada (28°C) & $P_{BOL,tot}$ & 1029,39 & kW & Antes de perdas; meta de barramento continua 1,0 MW \\
Espessura do Coverglass de Proteção & $t_{glass}$ & 100 & $\mu$m & Borosilicato ceriado CMG com revestimento MgF$_2$ \\
Escala do modelo DDD & $\Phi_x$ & $1{,}0 \times 10^{14}$ & e$^-$/cm$^2$ & Escolha paramétrica [PROVISÓRIA] \\
Coeficiente de ajuste pontual DDD & $C_d$ & 0,12851 & -- & Ajustado à retenção ECSS de 0,90 a $5\times10^{14}$ e$^-$/cm$^2$ [PROVISÓRIO] \\
Perda ambiental não radiativa (10 anos) & $\Delta P_{env}$ & 5,2 & \% & Premissa sem calibração ou ensaio [PROVISÓRIA] \\
\hline
\end{tabular}%
}
\end{table}

\begin{table}[!ht]
\centering
\footnotesize
\renewcommand{\arraystretch}{0.80}
\caption{Parâmetros eletroquímicos de 2ª ordem (DP) e constantes térmicas do BESS de 400~V.}
\label{tab:apendice_parametros_bess_dp}
\resizebox{0.95\textwidth}{!}{%
\begin{tabular}{lcccp{6.0cm}}
\hline
\textbf{Parâmetro / Propriedade} & \textbf{Símbolo} & \textbf{Valor Adotado} & \textbf{Unidade} & \textbf{Descrição e Referência Técnica} \\ \hline
Capacidade Nominal do Banco por Nó & $C_{n,no}$ & 22,52 & Ah & Pack 108S1P Li-ion NMC por satélite ($9{,}0\text{ kWh}$) \\
Energia Nominal do Banco por Nó & $E_{nom,no}$ & 9,0 & kWh & Energia instalada por satélite (60 kg) \\
Energia Utilizável Inicial ($DoD=60\%$) & $E_{util,no}$ & 5,4 & kWh & Margem segura para ciclo de vida $\ge 10$ anos LEO \\
Capacidade e Energia Total da Constelação & $C_{n,tot}, E_{tot}$ & 450,5 Ah / 180,0 kWh & -- & 20 nós satelitais operacionais interligados \\
Tensão Nominal do Pack & $V_{nom}$ & 399,6 & V & Platô de $3{,}70\text{ V/célula}$ a 50\% de SoC (108S) \\
Tensão Crítica de Subtensão & $V_{UV}$ & 324,0 & V & $3{,}0\text{ V/célula}$ (proteção contra dissolução do cobre) \\
Tensão Máxima de Recarga & $V_{OV}$ & 442,8 & V & $4{,}10\text{ V/célula}$ (proteção contra sobrecarga e gaseificação) \\
Resistência Ôhmica BOL (20°C) & $R_{0,ref}$ & 45,0 & m$\Omega$ & Impedância ôhmica pura do pack no início de vida \\
Resistência Ôhmica EOL (20°C) & $R_{0,eol}$ & 85,0 & m$\Omega$ & Degradação da camada SEI após 10 anos em LEO \\
Resistência de Transferência de Carga & $R_{p1}$ & 18,0 & m$\Omega$ & Cinética rápida de eletrodo / dupla camada elétrica \\
Capacitância de Dupla Camada & $C_{p1}$ & 3200,0 & F & Constante de tempo rápida $\tau_1 = R_{p1} C_{p1} = 57{,}6\text{ s}$ \\
Resistência de Difusão em Fase Sólida & $R_{p2}$ & 24,0 & m$\Omega$ & Transporte de massa de íons de lítio nos sítios ativos \\
Capacitância de Difusão & $C_{p2}$ & 22000,0 & F & Constante de tempo lenta $\tau_2 = R_{p2} C_{p2} = 528{,}0\text{ s}$ \\
Energia de Ativação de Arrhenius & $E_a$ & 0,32 & eV & Modulação térmica de transporte ($5{,}127 \times 10^{-20}\text{ J}$) \\
Massa Total do Pack de Baterias por Nó & $M_{bat}$ & 60,0 & kg & Densidade gravimétrica de $\approx 150\text{ Wh/kg}$ a nível de pack \\
Capacitância Térmica do Pack por Nó & $C_{th,bat}$ & 54360 & J/K & Massa e calor específico do NMC ($c_p \approx 906\text{ J/(kg K)}$) \\
Resistência Térmica da Manta MLI & $R_{th,MLI}$ & 4,8 & K/W & Isolamento passivo contra o chassi no vácuo \\
Potência Máxima do Aquecedor PWM & $P_{heat,max}$ & 300 & W & Aquecedor com modulação PWM (teto de segurança) \\
Rendimento Coulombico de Recarga & $\eta_c$ & 0,995 & -- & Eficiência de reversibilidade de carga \\
Dimensão do Vetor de Estado AUKF & $L$ & 3 & -- & $\mathbf{x} = [SoC, V_{p1}, V_{p2}]^T$ \\
Parâmetros de Espalhamento Sigma & $\alpha_{ut}, \beta_{ut}, \kappa_{ut}$ & $10^{-3}, 2{,}0, 0$ & -- & Parâmetros canônicos da Transformada Unscented \\
Covariância do Ruído de Processo & $\mathbf{Q}_0$ & $\text{diag}(10^{-7}, 10^{-5}, 10^{-5})$ & -- & Matriz de incerteza de transição de estado \\
Covariância do Ruído de Medição & $R_{meas}$ & 0,0225 & V$^2$ & Ruído gaussiano de sensoriamento ($150\text{ mV}_{rms}$) \\
Fator de Memória do Esquecimento & $\rho$ & 0,95 & -- & Filtro exponencial do resíduo de inovação AUKF \\
\hline
\end{tabular}%
}
\end{table}

\begin{table}[!ht]
\centering
\footnotesize
\renewcommand{\arraystretch}{0.80}
\caption{Parâmetros eletrofísicos, perdas e estabilidade dos estágios conversores CC-CC GaN.}
\label{tab:apendice_parametros_conversores}
\resizebox{0.95\textwidth}{!}{%
\begin{tabular}{lcccp{6.0cm}}
\hline
\textbf{Parâmetro / Propriedade} & \textbf{Símbolo} & \textbf{Valor Adotado} & \textbf{Unidade} & \textbf{Descrição e Referência Técnica} \\ \hline
Potência Nominal por Módulo Primário & $P_{mod}$ & 20,0 & kW & Módulo Full-Bridge 400 V / 48 V com GaN \\
Frequência de Chaveamento Primária & $f_{sw}$ & 150 & kHz & Otimização de volume magnético e perdas GaN \\
Resistência de Canal GaN ($25^\circ\text{C}$) & $R_{on,25}$ & 18,0 & m$\Omega$ & Transistor GaN HEMT 650 V classe espacial \\
Coeficiente Térmico do GaN & $\alpha_{sc}$ & 0,0065 & K$^{-1}$ & Inclinação positiva de $R_{on}(T_j)$ \\
Resistência de Canal GaN ($150^\circ\text{C}$) & $R_{on,150}$ & 32,63 & m$\Omega$ & Aumento de 81,25\% sob limite térmico máximo \\
Tempo de Subida do Transistor & $t_r$ & 12,0 & ns & Medição sob malha de gate otimizada de baixa indutância \\
Tempo de Descida do Transistor & $t_f$ & 15,0 & ns & Rotação de chaveamento com slew-rate de $33\text{ V/ns}$ \\
Capacitância Parasita Efetiva de Saída & $C_{oss}$ & 120,0 & pF & Capacitância equivalente dependente de tensão a 400 V \\
Carga de Porta e Tensão de Acionamento & $Q_g, V_{gate}$ & 6,0 / 5,0 & nC / V & Perdas de driver de $P_{gate} = 4{,}5\text{ mW}$ por dispositivo \\
Resistência Secundária Retificação ($25^\circ\text{C}$) & $R_{on,sec}$ & 0,80 & m$\Omega$ & Arranjo paralelo de chaves GaN 100 V a 48 V \\
Material Magnético do Transformador & -- & Ferrite MnZn & -- & Grau de alta frequência Ferroxcube 3C95 \\
Coeficiente de Steinmetz do Núcleo & $k_{fe}$ & 0,072 & W/(m$^3$Hz$^\alpha$T$^\beta$) & Caracterização volumétrica a $100^\circ\text{C}$ \\
Expoentes de Frequência e Indução & $\alpha, \beta$ & 1,46 / 2,52 & -- & Ajuste de curvas de Steinmetz MnZn \\
Volume Efetivo do Núcleo & $V_e$ & $4{,}5 \times 10^{-4}$ & m$^3$ & Geometria planar tipo E/PQ de alto acoplamento \\
Excursão de Densidade de Fluxo de Pico & $\Delta B$ & 0,15 & T & Ponto de operação abaixo da saturação ($B_{sat} \approx 0{,}4\text{ T}$) \\
Resistência CA dos Enrolamentos & $R_{ac,pri}, R_{ac,sec}$ & 15 / 0,25 & m$\Omega$ & Cabo Litz / fita de cobre com efeito pelicular a 150 kHz \\
Resistência Incremental CPL a 400 V & $r_{in,400}$ & $-8{,}00$ & $\Omega$ & Pequenos sinais $r_{in} = -V_{bus}^2/P$ para 20 kW \\
Resistência Incremental CPL a 48 V & $r_{in,48}$ & $-0{,}115$ & $\Omega$ & Pequenos sinais $r_{in} = -V_{bus}^2/P$ para 20 kW \\
Resistência Incremental CPL a 0,8 V & $r_{in,0.8}$ & $-0{,}00427$ & $\Omega$ & Carga local de acelerador de IA de 150 W \\
Capacitância de Barramento de Rack (48 V) & $C_{rack}$ & 15,0 & mF & Banco MLCC / tântalo com $ESR = 1{,}2\text{ m}\Omega$ nominal \\
Capacitância Degradada de Rack em EOL & $C_{rack,eol}$ & 3,0 & mF & Condição de pior caso WCA com $ESR = 8{,}0\text{ m}\Omega$ \\
\hline
\end{tabular}%
}
\end{table}

\vspace{0.8cm}
% ------------------------------------------------------------------------------
% APÊNDICE B
% ------------------------------------------------------------------------------
\subsection*{APÊNDICE B -- Deduções Algébricas e Modelagem Analítica Detalhada}
\addcontentsline{toc}{subsection}{APÊNDICE B -- Deduções Algébricas e Modelagem Analítica Detalhada}

\subsubsection*{B.1 Dedução da Relação de Massa entre Condutores de Alumínio e Cobre}

Para dois chicotes elétricos de mesmo comprimento $L$ e mesma resistência elétrica total $R$, a resistência ôhmica é dada por $R = \rho \frac{L}{A}$, onde $\rho$ é a resistividade e $A$ é a área de seção transversal condutora. Igualando as resistências dos condutores de Alumínio ($Al$) e Cobre ($Cu$):
\begin{equation}
    R_{Al} = R_{Cu} \implies \rho_{Al} \frac{L}{A_{Al}} = \rho_{Cu} \frac{L}{A_{Cu}} \implies \frac{A_{Al}}{A_{Cu}} = \frac{\rho_{Al}}{\rho_{Cu}}
\end{equation}

A massa total de cada condutor é expressa pelo produto do volume pela densidade mássica $\delta$:
\begin{equation}
    M_{Al} = \delta_{Al} \cdot A_{Al} \cdot L, \quad M_{Cu} = \delta_{Cu} \cdot A_{Cu} \cdot L
\end{equation}

Substituindo a razão de áreas na razão de massas:
\begin{equation}
    \frac{M_{Al}}{M_{Cu}} = \frac{\delta_{Al} \cdot A_{Al} \cdot L}{\delta_{Cu} \cdot A_{Cu} \cdot L} = \frac{\delta_{Al}}{\delta_{Cu}} \cdot \frac{\rho_{Al}}{\rho_{Cu}} = \frac{2{,}70\text{ g/cm}^3}{8{,}96\text{ g/cm}^3} \cdot \frac{2{,}82 \times 10^{-8}\ \Omega\cdot\text{m}}{1{,}72 \times 10^{-8}\ \Omega\cdot\text{m}} = 0{,}3013 \times 1,6395 \approx 0{,}494
\end{equation}
Demonstra-se analiticamente que o condutor de alumínio proporciona uma redução de massa de $1 - 0{,}494 = 0{,}506$ ($50{,}6\%$) em relação ao condutor de cobre equivalente.

\subsubsection*{B.2 Dedução da Corrente de Saturação Reversa da Célula Multijunção}

A equação do modelo de 1 diodo modificado é expressa por:
\begin{equation}
    I = I_{ph} - I_0 \left[ \exp\left( \frac{q (V + I R_s)}{a_{eq} k_B T} \right) - 1 \right] - \frac{V + I R_s}{R_{sh}}
\end{equation}

Na condição de circuito aberto ($V = V_{oc}$, $I = 0$):
\begin{equation}
    0 = I_{ph} - I_0 \left[ \exp\left( \frac{q V_{oc}}{a_{eq} k_B T} \right) - 1 \right] - \frac{V_{oc}}{R_{sh}}
\end{equation}

Isolando a corrente de saturação reversa $I_0$:
\begin{equation}
    I_0 \left[ \exp\left( \frac{q V_{oc}}{a_{eq} k_B T} \right) - 1 \right] = I_{ph} - \frac{V_{oc}}{R_{sh}} \implies I_0 = \frac{I_{ph} - \frac{V_{oc}}{R_{sh}}}{\exp\left( \frac{q V_{oc}}{a_{eq} k_B T} \right) - 1}
\end{equation}
Substituindo os valores de referência da Spectrolab XTJ Prime a $T_{ref} = 301{,}15\text{ K}$, com $V_{th} = 0{,}02595\text{ V}$ e $a_{eq} = 3{,}45$, obtém-se o valor unívoco de $I_{0,ref} = 3{,}4268 \times 10^{-14}\text{ A}$.

\subsubsection*{B.3 Formulação Matricial da Transformada Unscented e Fator de Esquecimento do AUKF}

Para um sistema não linear com dimensão de estado $L = 3$, o conjunto de $2L+1 = 7$ pontos sigma é calculado a partir da média $\hat{\mathbf{x}}_{k-1}$ e da matriz de covariância de erro $\mathbf{P}_{k-1}$ por:
\begin{align}
    \boldsymbol{\chi}_{0, k-1} &= \hat{\mathbf{x}}_{k-1} \\
    \boldsymbol{\chi}_{i, k-1} &= \hat{\mathbf{x}}_{k-1} + \left( \sqrt{(L + \lambda) \mathbf{P}_{k-1}} \right)_i, \quad i = 1, \dots, L \\
    \boldsymbol{\chi}_{i+L, k-1} &= \hat{\mathbf{x}}_{k-1} - \left( \sqrt{(L + \lambda) \mathbf{P}_{k-1}} \right)_i, \quad i = 1, \dots, L
\end{align}
onde o parâmetro composto de escala é $\lambda = \alpha_{ut}^2 (L + \kappa_{ut}) - L$. Os pesos de média $W_m$ e covariância $W_c$ são:
\begin{equation}
    W_m^{(0)} = \frac{\lambda}{L + \lambda}, \quad W_c^{(0)} = \frac{\lambda}{L + \lambda} + (1 - \alpha_{ut}^2 + \beta_{ut}), \quad W_m^{(i)} = W_c^{(i)} = \frac{1}{2(L + \lambda)}, \quad i = 1, \dots, 2L
\end{equation}

A matriz de raiz quadrada é calculada deterministicamente através da fatoração triangular inferior de Cholesky $\mathbf{S}\mathbf{S}^T = (L + \lambda)\mathbf{P}_{k-1}$. A propagação a priori dos estados obedece a $\boldsymbol{\chi}_{i, k|k-1} = f(\boldsymbol{\chi}_{i, k-1}, u_{k-1})$, resultando no estado previsto ponderado:
\begin{equation}
    \hat{\mathbf{x}}_{k|k-1} = \sum_{i=0}^{2L} W_m^{(i)} \boldsymbol{\chi}_{i, k|k-1}
\end{equation}

A covariância a priori $\mathbf{P}_{k|k-1}$ incorpora o fator de esquecimento adaptativo $\lambda_{ad,k} \ge 1$ para mitigar erros de modelagem sob transientes extremos de IA:
\begin{equation}
    \mathbf{P}_{k|k-1} = \lambda_{ad,k} \sum_{i=0}^{2L} W_c^{(i)} \left[ \boldsymbol{\chi}_{i, k|k-1} - \hat{\mathbf{x}}_{k|k-1} \right] \left[ \boldsymbol{\chi}_{i, k|k-1} - \hat{\mathbf{x}}_{k|k-1} \right]^T + \mathbf{Q}_k
\end{equation}
em que $\lambda_{ad,k}$ é modulado pela inovação $e_k = y_k - \hat{y}_{k|k-1}$ e pela covariância residual filtrada recursivamente $\mathcal{C}_{0,k} = \rho \mathcal{C}_{0,k-1} + (1 - \rho) e_k^2$ com memória $\rho = 0{,}95$:
\begin{equation}
    \lambda_{ad,k} = \max\left( 1, \, \frac{\mathcal{C}_{0,k} - R_{meas}}{P_{yy,k} - R_{meas}} \right)
\end{equation}
assegurando convergência estritamente estável mesmo sob desvios severos de condição inicial.

\subsubsection*{B.4 Condição de Fuga Térmica de GaN e Critério de Impedância de Middlebrook}

A condição de estabilidade térmica de regime estacionário para um semicondutor de potência operando com calor de condução termodependente $P_{loss}(T_j) = P_{cond}(T_j) + P_{sw}$ em contato térmico com dissipador em temperatura $T_c$ via resistência térmica $R_{th}$ é dada pelo balanço:
\begin{equation}
    f(T_j) = P_{loss}(T_j) - \frac{T_j - T_c}{R_{th}} = 0 \implies T_j = T_c + R_{th} P_{loss}(T_j)
\end{equation}
Diferenciando ambos os lados em relação a $T_j$:
\begin{equation}
    \frac{d f(T_j)}{d T_j} = \frac{\partial P_{loss}}{\partial T_j} - \frac{1}{R_{th}}
\end{equation}
Para que o ponto de equilíbrio seja assintoticamente estável frente a perturbações térmicas infinitesimais $\Delta T_j$, qualquer incremento de temperatura deve produzir uma taxa de remoção de calor estritamente superior ao incremento de geração interna, impondo:
\begin{equation}
    \frac{d f(T_j)}{d T_j} < 0 \implies \frac{\partial P_{loss}}{\partial T_j} < \frac{1}{R_{th}}
\end{equation}
Considerando $P_{cond}(T_j) = I_{rms}^2 R_{on,25} [1 + \alpha_{sc}(T_j - 25)]$ e perdas de comutação $P_{sw}$ aproximadamente invariantes com a temperatura, a derivada parcial reduz-se a:
\begin{equation}
    \frac{\partial P_{loss}}{\partial T_j} = I_{rms}^2 R_{on,25} \alpha_{sc} < \frac{1}{R_{th}} \implies R_{th} < \frac{1}{I_{rms}^2 R_{on,25} \alpha_{sc}}
\end{equation}
Caso $R_{th} \ge R_{th,crit}$, a curva linear de resfriamento tangencia ou cai abaixo da curva de geração de calor, resultando em ausência de raiz real estável e deflagrando a fuga térmica catastrófica (\textit{thermal runaway}).

No domínio da estabilidade de barramentos alimentando cargas de potência constante (CPL), a impedância de entrada incremental de pequenos sinais é obtida diferenciando a equação de potência constante $P_{CPL} = v(t) \cdot i(t)$:
\begin{equation}
    d P_{CPL} = 0 \implies i_0 \cdot \hat{v} + v_0 \cdot \hat{i} = 0 \implies Z_{in,CPL}(s) = \frac{\hat{v}(s)}{\hat{i}(s)} = - \frac{v_0}{i_0} = - \frac{V_{bus}^2}{P_{CPL}}
\end{equation}
A função de transferência de laço fechado da tensão de barramento sob interação fonte-carga é dada por:
\begin{equation}
    T_{bus}(s) = \frac{Z_{in,CPL}(s)}{Z_{out}(s) + Z_{in,CPL}(s)} = \frac{1}{1 + \frac{Z_{out}(s)}{Z_{in,CPL}(s)}} = \frac{1}{1 + T_m(s)}
\end{equation}
onde a razão de impedâncias minorante de Middlebrook é $T_m(s) = \frac{Z_{out}(s)}{Z_{in,CPL}(s)}$. Como $Z_{in,CPL}$ possui fase de $-180^\circ$ no plano complexo, o sistema em malha fechada possui polos no semiplano direito se o diagrama polar de $T_m(j\omega)$ envolver o ponto crítico $(-1, j0)$. Para evitar qualquer encadeamento oscilatório ou desamortecimento na presença de incertezas e variações de carga, estabelece-se a envoltória conservadora de Middlebrook:
\begin{equation}
    |Z_{out}(j\omega)| \le \frac{1}{GM} |Z_{in,CPL}(j\omega)| \implies |Z_{out}(j\omega)| \le 0{,}10 \frac{V_{bus}^2}{P_{CPL}}, \quad \forall \omega
\end{equation}

\subsubsection*{B.5 Dedução em Forma Fechada da Área de Sobreposição Cônica Solar}

Para a geometria de penumbra entre o disco solar de semi-ângulo aparente $\alpha_\odot$ e o disco terrestre de semi-ângulo aparente $\alpha_E$ separados pela distância angular de centros $\gamma$, a área de interseção de dois círculos $A_{ocultada}$ é deduzida pela soma dos dois setores circulares subtraída da área dos triângulos centrais correspondentes:
\begin{align}
    A_{ocultada}(\gamma) &= \alpha_\odot^2 \arccos\left( \frac{\gamma^2 + \alpha_\odot^2 - \alpha_E^2}{2 \gamma \alpha_\odot} \right) + \alpha_E^2 \arccos\left( \frac{\gamma^2 + \alpha_E^2 - \alpha_\odot^2}{2 \gamma \alpha_E} \right) \nonumber \\
    &\quad - \frac{1}{2} \sqrt{\left[ (\alpha_\odot + \alpha_E)^2 - \gamma^2 \right] \left[ \gamma^2 - (\alpha_E - \alpha_\odot)^2 \right]}
\end{align}
Essa expressão fornece transição estritamente contínua de classe $C^1$ entre iluminação plena ($S=1$) e umbra total ($S=0$).

\subsubsection*{B.6 Formulação Integral do Fator de Visão de Calota Esférica (ECSS)}

O fator de visão geométrico $F_{k-E}$ para uma superfície plana com inclinação $\phi_k$ em relação ao vetor Nadir a partir de uma calota esférica visível com semi-ângulo $\theta_m = \arcsin(1/H)$, onde $H = (R_E + h)/R_E$, obedece à integração analítica da norma ECSS-E-HB-31-01 Parte 1:
\begin{equation}
    F_{k-E}(h, \phi_k) = 
    \begin{cases}
        \frac{1}{H^2} \cos \phi_k, & \text{se } 0 \le \phi_k \le \frac{\pi}{2} - \theta_m \\[0.15cm]
        \frac{1}{2} - \frac{1}{\pi} \arcsin\left( \frac{\sqrt{H^2 - 1}}{H \sin \phi_k} \right) + \frac{\cos \phi_k}{\pi H^2} \arccos\left( \sqrt{H^2 - 1} \cot \phi_k \right) - \frac{\sqrt{H^2 - 1}}{\pi H^2} \sqrt{1 - H^2 \cos^2 \phi_k}, & \text{se } |\phi_k - \frac{\pi}{2}| \le \theta_m \\[0.15cm]
        0{,}0, & \text{se } \phi_k > \frac{\pi}{2} + \theta_m
    \end{cases}
\end{equation}

\subsubsection*{B.7 Modelo BSIM de Corrente de Fuga Sublimiar de Transistores Nanométricos}

A dependência termodinâmica da corrente de fuga por tunelamento e emissão termoiônica em transistores FinFET e GAA sub-7~nm é expressa pela formulação de referência BSIM:
\begin{equation}
    P_{leak}(T_j, V_{core}) = V_{core} I_{leak,0}(V_{core}) \left( \frac{T_j}{T_{ref}} \right)^2 \exp\left[ \frac{-q \left( V_{th0} - \kappa_T (T_j - T_{ref}) \right)}{\eta_{sub} k_B T_j} \right] \left[ 1 - \exp\left( -\frac{q V_{core}}{k_B T_j} \right) \right]
\end{equation}
com parâmetros $V_{th0} \approx 0{,}28\text{ V}$, $\kappa_T \approx 1{,}10\text{ mV/K}$ e $\eta_{sub} \approx 1{,}25$.

\vspace{0.8cm}
% ------------------------------------------------------------------------------
% APÊNDICE C
% ------------------------------------------------------------------------------
\subsection*{APÊNDICE C -- Rotinas Computacionais em GNU Octave}
\addcontentsline{toc}{subsection}{APÊNDICE C -- Rotinas Computacionais em GNU Octave}

As rotinas computacionais desenvolvidas para as simulações multifísicas encontram-se estruturadas no diretório \texttt{simulacoes/} do repositório da pesquisa:
\begin{itemize}[noitemsep,topsep=2pt,leftmargin=*]
    \item \texttt{simulacao\_3\_1\_barramento\_hvdc.m}: Dinâmica do barramento de 400~V, balanço orbital multicíclico do BESS (5 órbitas), amortecimento RLC e estresse $I^2t$ do SSPC.
    \item \texttt{simulacao\_3\_2\_gerador\_fotovoltaico\_am0.m}: Curvas $I-V/P-V$ multijunção, degradação temporal DDD (10 anos), avalanche reversa em hot spot e dinâmica do MPPT.
    \item \texttt{gerar\_figuras\_compostas\_3\_1.m}: Consolidação de dados e renderização das Figuras 3.1 (balanço BESS) e 3.2 (curto-circuito e acoplamento térmico LHP).
    \item \texttt{gerar\_figuras\_compostas\_3\_2.m}: Renderização das Figuras 3.3 (curvas AM0 térmicas e DDD) e 3.4 (hot spot e transiente do MPPT GaN).
    \item \texttt{gerar\_figuras\_compostas\_3\_3.m}: Modelo eletroquímico Dual-Polarization (DP) 2ª ordem, cinética de Arrhenius, estimador AUKF e colapso por subtensão (Figuras 3.5 e 3.6).
    \item \texttt{gerar\_figuras\_compostas\_3\_4.m}: Modelo eletrotérmico GaN, perdas de Steinmetz a 150~kHz, estabilidade de Middlebrook sob CPL e desacoplamento de rack a 48~V (Figuras 3.7 e 3.8).
    \item \texttt{simulacao\_4\_1\_fluxos\_orbitais.m}: Propagação kepleriana, ocultação cônica exata, fatores de forma ECSS $F_{rad-E}(h, \phi)$, fluxos ambientais e sensibilidade de atitude.
    \item \texttt{simulacao\_4\_6\_fadiga\_bga\_pue.m}: Modelo Coffin-Manson-Engelmaier para soldas BGA SAC305, matriz fatorial orbital/TEMS, indicador $PUE_{orbital}$ e rendimento computacional TOPS/W.
    \item \texttt{gerar\_figuras\_compostas\_capitulo4.py}: Script multifísico para renderização das Figuras 4.1 (fluxos orbitais), 4.2 (rede nodal 5 nós), 4.3 (co-design FinFET/runaway) e 4.4 (NMPC).
\end{itemize}

